\section{Week 10}

\subsection{Corollary to Orbit-Stabilzer Theorem} 
\begin{definition}[ ]
    Let \( G  \) act on \( A  \) where the action is \( g \cdot a  \)
    \begin{enumerate}
        \item[(1)] The equivalence class \( \{  g \cdot a : g \in G  \} \subseteq A   \) is called the orbit of \( G  \) containing \( a \in A  \). We denote \( \mathcal{O}_{a} \).
        \item[(2)] The action of \( G  \) on \( A  \) is called transitive if there is only one orbit; that is, given \( a,b \in A  \), there exists a \( g \in G  \) such that \( a = g \cdot b  \).
    \end{enumerate}
\end{definition}

\begin{eg}
   Let \( G = {S}_{n} \) act on \( A = \{  1,2, \dots, n \}  \) in the usual way. Then observe that  
   \begin{align*}
       {\mathcal{O}}_{2} &= \{  g \cdot 2 : g \in {S}_{n}  \}  \\
       &= \{  \sigma \cdot 2 : \sigma \in {S}_{n} \}  \\
       &= \{  \sigma(2) : \sigma \in {S}_{n} \} \\
       &= \{ 1,2,3,\dots n \}  = A,
    \end{align*}
    but this orbit does not really give us any information on the structure of \( A  \).
\end{eg}

\begin{eg}
   Let \( \sigma \in S \). Let's say that  
   \begin{align*}
       \sigma(1) &= 2 \\
       \sigma(2) &= 3 \\
       \sigma(3) &= 4 \\
       \sigma(4) &= 1 \\
       \sigma(5) &= 6 \\
       \sigma(6) &= 5
   \end{align*}
   and so we have 
   \[  \sigma = (1 \ 2 \ 3 \ 4 )(5 \ 6) \]
   Let \( G = \langle \sigma  \rangle \leq {S}_{7}  \). By the definition of a group action, if a group acts on a set \( A  \), then any subgroup of the group acts on \( A  \). This is the case because \( g \cdot a  \) is a mapping from \( G \times A  \) to \( A \). So, \( g \cdot a  \) is an element of \( A \) for all \( g \in G  \) and \( a \in A  \). Thus, we have an action when we restrict to \( g  \) to a subgroup. Let \( \langle  \sigma \rangle \) act on \( A = \{  1,2,3,\dots, 7 \}   \) in the usual way. Next, we will compute the orbits relating to this action. That is, for \( | \langle \sigma  \rangle |  = 3 \)
   \begin{align*}
       {\mathcal{O}}_{1} &= \{  \pi(1) : \pi \in \langle \sigma  \rangle \} \\
                         &= \{ \sigma^{k}(1) : k \in \{ 0,1,2,3 \}  \}  \\
                         &= \{ 1 , \sigma(1), \sigma^{2}(1), \sigma^{3}(1) \}  \\
                         &= \{ 1,2,3,4 \} \\
       {\mathcal{O}}_{5} &= \{  5 , \sigma(5), \sigma^{2}(5), \sigma^{3}(5) \}  \\
                         &= \{ 5,6 \} \\
                         &= {\mathcal{O}}_{6} \\
        {\mathcal{O}}_{7} &= \{ 7 \}.
   \end{align*}
\end{eg}

\begin{corollary}[To Orbit Stabilizer Theorem]
    This tells us that every \( \sigma \in {S}_{n} \) can be written as the product of disjoint cycles
\end{corollary}
\begin{proof}
Let \( \sigma \in {S}_{n} \) and \( A = \{  1,2,\dots, n \}  \). Also, let \( G = \langle  \sigma \rangle \). Then \( G  \) acts on \( A  \) in the usual way. Then by the Orbit-Stabilzer Theorem, this action partitions \( A  \) into a unique set of disjoint orbits. Let \( {\mathcal{O}}_{x}  =  \{  \sigma^{k} : k \in \Z \}  \) be one such orbit where \( x \in A \). Again, by the Orbit Stabilizer Theorem applied to \( {\mathcal{O}}_{x}  \) there is a bijection between \( {\mathcal{O}}_{x}  \) and \( G / {G}_{x} \) where \( {G}_{x} = \langle \sigma   \rangle / {G}_{x} \) is the stabilizer with respect to the given action. This bijection is defined by \( \sigma^{k}(x) \mapsto \sigma^{k} {G}_{x} \). Since \( G  \) is cyclic, it is abelian and so \( {G}_{x} \trianglelefteq G  \). Hence, \( G / {G}_{x} \) is cyclic of order \( d  \) where \( d  \) is the smallest positive integer such that \( \sigma^{d} \in {G}_{x} \). Also, \( d = | G : {G}_{x} | = | {\mathcal{O}}_{x} |    \). This tells us the different cosets of \( {G}_{x} \) in \( G = \langle \sigma  \rangle \) are   
\[  1 {G}_{x}, \sigma {G}_{x} , \sigma^{2} {G}_{x}, \dots, \sigma^{d-1} {G}_{x} \]
and the corresponding elements of \( {\mathcal{O}}_{x}  \) are \( x, \sigma(x) , \sigma^{2}(x), \dots, \sigma^{d-1}(x) ,  (\sigma^{2}(x)) \). Thus, \( \sigma \) cycles the elements of \( {\mathcal{O}}_{x} \); that is, the elements are a cycle size \( d  \). Hence, \( \sigma \) creates a \( d-\) cycle. Since the orbits form a partition of \( A  \) each orbit gives a disjoint cycle in our decomposition of \( \sigma \).
\end{proof}

\subsection{Groups Acting on Themselves by Left-Multiplication}

Let \( G  \) be a group, let \( A = G  \), and define \( (\cdot) : G \times A \to G  \) by \( g \cdot a = ga  \) for all \( g \in G  \) and \( a \in G  \). Let \( \psi  \) be the \textbf{permutation representation of this action}. Recall that \( \psi : G \to {S}_{A} \) where \( A = G  \) and is defined by \( \psi(g) = {\sigma}_{g} \in {S}_{A} \) where \( {\sigma}_{g}(a) = g \cdot a  \). 

\begin{eg}
   Let \( G = {V}_{4} = \{  1,a,b,c \}  \) where \( a^{2} = b^{2} = c^{2} = 1  \). Let's see the permutation representation in action when \( G  \) acts on itself by Left-multiplication. Let's look at the following mapping:
   \begin{align*}
       {\sigma}_{1}(x) &= {\sigma}_{1}(1) = 1 \cdot 1 = {1}_{{S}_{4}} \implies {\sigma}_{1}(a) = 1 \cdot a = a    \\
       {\sigma}_{1}(b) \ \ &\text{and} \ \  {\sigma}_{1}(c) = c 
   \end{align*}
   Also, \( {\sigma}_{a}(1) = a 1  =a  \), \( {\sigma}_{a}(a) = 1  \), \( {\sigma}_{a}(b) = c  \), \( {\sigma}_{a}(c) = b  \), and so 
   \[  {\sigma}_{a} = (1 \ a )(b \ c). \]
   Also, \( {\sigma}_{b}(1) = b1 = b  \), and \( {\sigma}_{b}(a) = c   \), and, \( {\sigma}_{c}(b) = a  \), and so   
   \[  {\sigma}_{b} = (1 \ b)(a \ c) \]
   and 
   \[  {\sigma}_{c} = (1 \ c)(a \ b). \]
   We know that \( \psi : {V}_{4} \to {S}_{4} \) and so by the first isomorphism theorem, we can say that 
   \[  {V}_{4} / \text{ker}(\psi) \cong \psi({V}_{4}) = \{ {1}_{{S}_{4}} , (1 \ 2)(3 \ 4), (1 \ 3)(2 \ 4), (1 \ 4)(2 \ 3) \}  \]
   and
   \[ \text{ker}(\varphi) = \{ g \in G : g \cdot a  = a \ \   \forall a \in A   \} = \{ {1}_{{V}_{4}} \}.    \]
   Thus, we can say that \( {V}_{4} \) is isomorphic to a subgroup of \( {S}_{4} \); that is, \( {V}_{4} \cong \psi({V}_{4}) \).
\end{eg}
\subsection{The generalization of \( G  \) acting on itself by Left-multiplication.}

 Let \( G  \) be a group and let \(  H \leq G  \). Let \( A = G / H  \). We'll define the action on \( G  \) on \( A  \) by \( g \cdot aH  = (ga)H \) and this is for all \(  g \in G  \) and \( a H \in A  \).
 Note if \( H = \{ {1}_{G} \}  \), then \( g \cdot aH = ga H  = ga \{ {1}_{G} \}  = \{ ga \}  \).

 \begin{eg}
    Let \( G = {D}_{8} \) and \( H = \langle s  \rangle = \{  1,s \}   \). The index of \( \langle s  \rangle \) in \( {D}_{8} \) is 4. The cosets of \( H  \) are \( 1 H  \), \( r H  \), \( r^{2} H  \), and \( r^{3} H  \) where let's label them 1-4 accordingly. Lets find the permutation representation of this action. We don't have to work too hard since we know that \( r  \) and \( s  \) are the generators of \( {D}_{8} \). That is let us just consider \( {\sigma}_{s} \) and \( {\sigma}_{r} \). Then \( s \cdot 1 H = \)   
 \end{eg}


 
